\chapter{正比例函数与一次函数}\label{ch:g9:linfunc}

“三千克的价钱是一千克的三倍”，这正是成正比例，
也是两个量之间最简单的关系。正比例函数是用函数语言表达的成正比例关系；
一次函数则在正比例函数的基础上，加上一个固定的起始量。
它们的图像都是直线，读懂这些直线
正是本章要培养的技能。

\section{成正比例与正比例函数}

\begin{definition}[正比例函数]\label{def:g9:linfunc:linear}
称\emph{正比例函数}\index{正比例函数}为把每个数乘以一个
固定的数 $a$：
\[
f(x) = a x .
\]
数 $a$ 称为 $f$ 的\emph{系数}。两个量
\emph{成正比例}\index{成正比例}，当且仅当其中一个量是
另一个量的正比例函数。
\end{definition}

\begin{example}\label{ex:g9:linfunc:linear}
若苹果每千克 $3$ 元，则 $x$ 千克苹果的价格为
$f(x) = 3x$：购买两倍的数量，花费也是两倍；
$2.5$ 千克的价格为 $f(2.5) = 7.5$。

反之，若正比例函数满足 $f(4) = 10$，则
$a = \frac{10}{4} = 2.5$，$f(x) = 2.5x$：\emph{一个函数值
就能完全确定一个正比例函数。}
\end{example}

\begin{proposition}[正比例函数的图像]\label{prop:g9:linfunc:lineargraph}
$f(x) = ax$ 的图像是一条\emph{经过原点}的直线。
反之，经过原点的每条非竖直直线都是
某个正比例函数的图像。
\end{proposition}
\admitted

\begin{omfigure}
\begin{tikzpicture}
\begin{axis}[omaxis,
  width=7.4cm, height=5.6cm,
  xmin=-3.4, xmax=3.4, ymin=-3.2, ymax=3.6,
  xtick={-2,-1,1,2}, ytick={-2,-1,1,2}]
  \addplot[omDef, thick, domain=-2.2:2.2] {1.5*x};
  \addplot[omThm, thick, domain=-3.1:3.1] {-0.5*x};
  \node[omDef, font=\small] at (axis cs:2.15,2.6) {$y = 1.5x$};
  \node[omThm, font=\small] at (axis cs:-2.5,1.7) {$y = -0.5x$};
\end{axis}
\end{tikzpicture}

{\small 两个正比例函数：两者的图像都经过原点；
系数 $a$ 是 $x = 1$ 时图像达到的高度。}
\end{omfigure}

\section{一次函数}

\begin{definition}[一次函数]\label{def:g9:linfunc:affine}
称\emph{一次函数}\index{一次函数}具有如下形式：
\[
f(x) = a x + b ,
\]
其中 $a$ 与 $b$ 为固定的数。正比例函数是
$b = 0$ 时的特殊情形；常值函数则是 $a = 0$ 的情形。
\end{definition}

\begin{example}
某健身房收取 $20$ 元报名费，此后每次健身收费 $5$ 元，
则 $x$ 次健身的总费用为 $f(x) = 5x + 20$。费用\emph{不}
与次数成正比例（10 次的费用不是 5 次的两倍），
但每多一次，费用都增加相同的数额 $5$。
\end{example}

\begin{proposition}[一次函数的图像]\label{prop:g9:linfunc:affinegraph}
$f(x) = ax + b$ 的图像是一条直线：
\begin{itemize}
  \item $b = f(0)$ 是\emph{$y$ 轴截距}：直线
        与 y 轴相交于点 $(0, b)$；
  \item $a$ 是\emph{斜率}\index{斜率}：向右移动一个单位，
        对应上升 $a$ 个单位（$a < 0$ 时为下降）；函数
        在 $a > 0$ 时递增，在 $a < 0$ 时递减；
  \item 对任意两个输入 $u \neq v$，
        $a = \dfrac{f(v) - f(u)}{v - u}$.
\end{itemize}
\end{proposition}

\begin{proof}[最后一条的证明]
$f(v) - f(u) = (av + b) - (au + b) = a(v - u)$，两边除以 $v - u$。
其余部分在本阶段不加证明地承认（高中分册对直线及其方程给出了
完整的论述）。
\end{proof}

\begin{omfigure}
\begin{tikzpicture}
\begin{axis}[omaxis,
  width=7.8cm, height=5.8cm,
  xmin=-1.9, xmax=4.4, ymin=-1.4, ymax=5.4,
  xtick={1,2,3}, ytick={1,2,3}]
  \addplot[omDef, thick, domain=-1.6:4.1] {x + 1};
  \fill (axis cs:0,1) circle (1.5pt);
  \node[left, font=\small] at (axis cs:0,1.25) {$b$};
  \draw[dashed, omThm] (axis cs:1,2) -- (axis cs:2,2) -- (axis cs:2,3);
  \node[font=\footnotesize, below] at (axis cs:1.5,2) {$+1$};
  \node[font=\footnotesize, right] at (axis cs:2,2.5) {$+a$};
\end{axis}
\end{tikzpicture}

{\small 从图像读出 $y = ax + b$：直线与 $y$ 轴
相交于高度 $b$ 处，向右每移动一个单位便上升 $a$（此处
$a = 1$，$b = 1$）。}
\end{omfigure}

\begin{method}[由两个函数值求一次函数]\label{met:g9:linfunc:findaffine}
已知 $f(u)$ 与 $f(v)$，其中 $u \neq v$：
\begin{enumerate}
  \item 计算斜率 $a = \dfrac{f(v) - f(u)}{v - u}$；
  \item 将两个已知值之一代入 $f(x) = ax + b$，
        求出 $b$；
  \item 用另一个值进行检验。
\end{enumerate}
\end{method}

\begin{example}\label{ex:g9:linfunc:findaffine}
求满足 $f(2) = 7$ 且 $f(5) = 16$ 的一次函数。
\[
a = \frac{16 - 7}{5 - 2} = \frac93 = 3 ,
\]
于是由 $f(2) = 3 \times 2 + b = 7$ 得 $b = 1$，$f(x) = 3x + 1$。检验：
$f(5) = 15 + 1 = 16$。
\end{example}

\section{百分数}

\begin{proposition}[百分比变化]\label{prop:g9:linfunc:percent}
将一个量增加 $t\,\%$，则它乘以
$1 + \dfrac{t}{100}$；将其减少 $t\,\%$，则它乘以
$1 - \dfrac{t}{100}$。百分比变化都是正比例函数。
\end{proposition}

\begin{proof}
将 $x$ 增加 $t\,\%$，即加上 $\frac{t}{100}\,x$：
\[
x + \frac{t}{100}\,x = \left(1 + \frac{t}{100}\right) x .
\]
减少时计算相同，只需将加号改为减号。
\end{proof}

\begin{example}\label{ex:g9:linfunc:percent}
价格为 $80$ 的商品上涨 $15\%$，新价格为
$80 \times 1.15 = 92$。价格为 $60$ 的商品降价 $30\%$：
$60 \times 0.7 = 42$。

\emph{连续变化时，倍数相乘。}先上涨 $20\%$，
再下降 $20\%$，得到
$1.2 \times 0.8 = 0.96$：总体上是 $4\%$ 的 \emph{下降}，而不是
回到原价！
\end{example}

\section{练习}

\begin{exercise}[$\star$]\label{exo:g9:linfunc:1}
判断下列函数中，哪些是正比例函数，哪些是一次函数，哪些两者都不是。
\[
f(x) = 4x, \qquad
g(x) = 3x - 5, \qquad
h(x) = x^2, \qquad
k(x) = -\tfrac x2, \qquad
m(x) = 7 .
\]
\end{exercise}

\begin{exercise}[$\star$]\label{exo:g9:linfunc:2}
设 $f(x) = 2x - 3$。计算 $f(0)$、$f(4)$、$f(-1)$，并求出数
$x$，使 $f(x) = 9$。
\end{exercise}

\begin{exercise}[$\star$]\label{exo:g9:linfunc:3}
一个正比例函数满足 $f(6) = 15$。求其系数，再
计算 $f(10)$ 与 $f(-2)$。
\end{exercise}

\begin{exercise}[$\star$]\label{exo:g9:linfunc:4}
在同一坐标系中画出 $f(x) = 2x$、
$g(x) = 2x + 3$ 与 $h(x) = -x + 4$ 的图像。关于
$f$ 与 $g$ 的图像，能得出什么结论？
\end{exercise}

\begin{exercise}[$\star$]\label{exo:g9:linfunc:5}
求满足 $f(1) = 5$ 且 $f(3) = 11$ 的一次函数；再求满足
$g(0) = 4$ 且 $g(2) = 0$ 的一次函数。
\end{exercise}

\begin{exercise}[$\star\star$]\label{exo:g9:linfunc:6}
计算：上涨 $25\%$ 后 $120$ 元的价格；
降价 $40\%$ 后 $65$ 元的价格；一件商品
售价为 $69$（已降价 $25\%$），求原价。
\end{exercise}

\begin{exercise}[$\star\star$]\label{exo:g9:linfunc:7}
某手机套餐 A 每月收费 $10$ 元，另加每分钟通话 $0.05$ 元；套餐
B 每月收费 $25$ 元，通话时长不限。
\begin{enumerate}
  \item 将套餐 A 的月费表示为通话分钟数
        $x$ 的函数。
  \item 每月通话超过多少分钟时，套餐 B 更便宜？
\end{enumerate}
\end{exercise}

\begin{exercise}[$\star\star$]\label{exo:g9:linfunc:8}
某细菌群体初始有 $2000$ 个细菌，每小时增长 $10\%$。
\begin{enumerate}
  \item 一小时后与两小时后各有多少个细菌？
  \item 说明为什么两小时后的个数不是 $2400$。
\end{enumerate}
\end{exercise}

\begin{exercise}[$\star\star$]\label{exo:g9:linfunc:9}
某一次函数的图像经过点 $(2, 3)$ 与
$(6, 1)$。
\begin{enumerate}
  \item 求其斜率，该函数是单调递增还是单调递减？
  \item 求其表达式，并求使函数值
        为 $0$ 的自变量的值。
\end{enumerate}
\end{exercise}

\begin{exercise}[$\star\star\star$]\label{exo:g9:linfunc:10}
某商店先将价格提高 $t\,\%$，再将新价格降低
$t\,\%$。
\begin{enumerate}
  \item 证明：最终价格等于原价乘以
        $1 - \dfrac{t^2}{10000}$.
  \item 由此推出最终价格总是\emph{低于}
        原价（当 $0 < t \leq 100$ 时），并求
        $t = 10$ 时的总体百分比变化。
\end{enumerate}
\end{exercise}

\section{问题：两支温度计}

\begin{problem}[{周末问题 --- 摄氏度、华氏度与生活中的一次函数：某一温度
在两种刻度上读数相同，而恒定斜率正是
直线性的标志}]
\label{pb:g9:linfunc:1}
美国人听到“$68$ 度”便穿得很随意；欧洲人听到
这个数字，便赶紧拿上外套。世界上的两支温度计以一次函数相联系——
由两个事实构造它，正是
\cref{met:g9:linfunc:findaffine}所体现的技能。本题
构造换算公式，求出两种刻度读数一致的那个奇特温度，
随后用同一副眼镜读出出租车、蟋蟀与蜡烛的规律，
并以一项能够识别直线性本身的检验
作结。

\medskip
\noindent\textbf{第一部分 --- 构造换算公式。}
两个事实确定了华氏温标：水在
$32\,^\circ$F 结冰，在 $212\,^\circ$F 沸腾（对应 $0\,^\circ$C 与
$100\,^\circ$C）。
\begin{enumerate}
  \item 求将摄氏度换算为华氏度的一次函数 $F = f(C)$
        （用\cref{met:g9:linfunc:findaffine}，
        并取 $f(0) = 32$ 与 $f(100) = 212$）。
  \item 换算成华氏度：温和的 $20\,^\circ$C 一天；
        人体的 $37\,^\circ$C；寒冷的 $-10\,^\circ$C。
  \item 解出 $C$ 以建立反向换算，并用它换算
        $50\,^\circ$F 与 $98.6\,^\circ$F。
  \item $f$ 是否为\emph{正比例函数}
        （见\cref{def:g9:linfunc:linear}）？检验其判别标志：
        $f(0)$ 是多少？摄氏读数加倍时，
        华氏读数是否也加倍？据此给出正确的名称
        （见\cref{def:g9:linfunc:affine}）。
  \item 经典谜题：是否存在这样一个温度，它在两种刻度上读出
        \emph{同一个数}？解方程 $f(x) = x$，
        并说明该解在 $f$ 的图像上意味着什么
        （图像在该处与哪条直线相交？）。
\end{enumerate}

\medskip
\noindent\textbf{第二部分 --- 生活中的一次函数。}
\begin{enumerate}[resume]
  \item 出租车 A 收取 $3$ 欧元，另加每千米 $1.50$ 欧元；出租车 B
        统一按每千米 $2$ 欧元收费。写出两种价格函数，
        并比较 $4$ 千米与
        $8$ 千米行程的费用。
  \item 求两车费用相同时的里程，
        并在图像上描述这一情形（两条直线 ——
        交点处发生了什么？交点之前与之后分别如何？
        ）。
  \item 博物学家的经验法则：蟋蟀每分钟约鸣叫
        $N = 7T - 30$ 次（温度为
        $T\,^\circ$C 时）。在 $20\,^\circ$C 的
        傍晚，蟋蟀会鸣叫多少次？若你一分钟内数到 $82$ 次，
        蟋蟀所报出的温度是多少？
  \item 一部以 $600$ 欧元购入的手机，每月贬值
        $15$ 欧元：写出价值函数 $v(t)$，求
        手机价值为 $240$ 欧元的时刻，以及模型表明它
        价值为零的时刻。
        该日期之后，这个公式是否仍然有意义？
  \item 百分比变化即正比例函数：$+25\,\%$ 相当于
        乘以 $1.25$，而 $-20\,\%$ 相当于乘以 $0.8$。
        （见\cref{prop:g9:linfunc:percent}）。将它们复合：
        先上涨 $+25\,\%$，再作 $-20\,\%$ 的调整，对一个价格有何影响？
        这一小小的“奇迹”何以如此？请解释，并与
        \cref{exo:g9:linfunc:10}作比较。
\end{enumerate}

\medskip
\noindent\textbf{第三部分 --- 斜率决定一切。}
\begin{enumerate}[resume]
  \item 一支蜡烛点燃时高 $20$ 厘米，半小时后
        高 $17$ 厘米。假设燃烧呈一次函数规律，求 $h(t)$
        （高度单位为厘米，$t$ 单位为分钟），并预测蜡烛何时燃尽。
  \item 两名骑行者沿同一条路骑行：第一人的路程为
        $d_1(t) = 24t$（$t$ 小时后的路程，单位为千米），第二人的为
        $d_2(t) = 10 + 18t$。谁具有起步优势，
        谁骑得更快，第一人在何时、何处
        追上第二人？
  \item 设 $f(x) = 2x + 3$，$g(x) = -0.5x + 7$。计算
        $f(x) + g(x)$，并验证它仍是一次函数。一般地，两个一次函数之和的
        斜率与截距各是什么？
  \item 完整讨论方程 $mx + p = 0$：
        当 $m \neq 0$ 时有多少个解（并求出该解）？当 $m = 0$
        且 $p \neq 0$ 时如何？当 $m = 0$ 且 $p = 0$ 时如何？（每种情形
        都有图像上的对应：斜率为 $m$ 的直线与
        水平轴在何处相交？）
  \item 共线检验：点 $(1, 5)$、$(3, 9)$、
        $(7, 17)$ 是否在同一条直线上？点 $(1, 5)$、$(3, 9)$、
        $(6, 16)$ 呢？计算各对点之间的
        变化率并作出判断。总结：\emph{恒定的
        变化率}是一次函数的标志——曲线图像会改变其
        变化率，而度量\emph{这一变化}则属于
        高中分册的内容。
\end{enumerate}
\end{problem}
